% Generated draft: requires mathematical and source review.
\documentclass[11pt]{article}
\usepackage[margin=25mm]{geometry}
\usepackage{fontspec}
\setmainfont{Latin Modern Roman}
\usepackage{amsmath,amssymb,mathtools,unicode-math}
\setmathfont{Latin Modern Math}
\usepackage{enumitem,needspace,xurl,hyperref}
\setlength{\parindent}{0pt}
\setlength{\parskip}{5pt}
\setlength{\emergencystretch}{3em}
\newcommand{\N}{\mathbb{N}}
\newcommand{\Z}{\mathbb{Z}}
\newcommand{\Q}{\mathbb{Q}}
\newcommand{\R}{\mathbb{R}}
\newcommand{\C}{\mathbb{C}}
\newcommand{\F}{\mathbb{F}}
\newcommand{\A}{\mathbb{A}}
\newcommand{\PP}{\mathbb{P}}
\newcommand{\E}{\mathbb{E}}
\newcommand{\eps}{\varepsilon}
\newcommand{\dd}{\,\mathrm{d}}
\begin{document}
\section*{UnsolvedMath ID: 801116}
\textbf{Source problem code:} catalogue entry
\subsection{Definitions and mathematical statement}
Let \(\mathbb F_2\) denote the free group on two generators. For a discrete group \(G\), let \(\ell^2(G)\) have its canonical orthonormal basis \((\delta_h)_{h\in G}\), and define the left regular representation by \(\lambda_G(g)\delta_h=\delta_{gh}\). Its group von Neumann algebra and canonical trace are
\[
L(G)=\{\lambda_G(g):g\in G\}'',\qquad \tau_G(x)=\langle x\delta_e,\delta_e\rangle,
\]
where the double commutant is taken in \(B(\ell^2(G))\) and \(e\) is the identity of \(G\).

A group is \emph{icc} if it is infinite and every conjugacy class of a nonidentity element is infinite. It is \emph{inner amenable} if there exists a finitely additive probability measure \(m\) on all subsets of \(G\setminus\{e\}\) such that
\[
m(gAg^{-1})=m(A)\qquad\text{for every }g\in G\text{ and }A\subseteq G\setminus\{e\}.
\]
Thus non-inner-amenability means that no such measure exists.

Under the standard interpretation in the continuous first-order language of tracial von Neumann algebras, write \((M,\tau_M)\equiv(N,\tau_N)\) when every sentence has the same value in the two structures. The source question is:
\[
\text{Does there exist a discrete icc group }G\text{ that is not inner amenable and satisfies }
(L(G),\tau_G)\not\equiv(L(\mathbb F_2),\tau_{\mathbb F_2})\,?
\]
Equivalently, the last condition asks for a sentence \(\varphi\) in that language with \(\varphi^{(L(G),\tau_G)}\ne\varphi^{(L(\mathbb F_2),\tau_{\mathbb F_2})}\).

The discrete-group setting and tracial logical language above make explicit the standard interpretation of the notation; the supplied statement does not specify the language or whether \(G\) must be countable. No countability restriction is imposed in this draft, and the intended restriction requires review. The supplied statement carries a solved annotation attributing a solution to J. Peterson, whereas the registry status is \texttt{open}. These are conflicting supplied labels, not an independently verified assessment.
\subsection{Short English statement}
Does there exist an infinite discrete group whose nonidentity conjugacy classes are all infinite, which has no conjugation-invariant mean on its nonidentity elements, and whose group von Neumann algebra has a different first-order tracial theory from that of the free group on two generators? Countability and the precise logical language require confirmation. The source's solved annotation conflicts with the registry's open label.
\subsection{Sources}
\begin{itemize}
\item Local UnsolvedMath registry, permanent ID \(801116\), source problem code \texttt{catalogue entry}, title: \emph{Ultraproduct \(\mathrm{II}_1\) Factors --- Problem 17}. The supplied record attributes the problem to Elayavalli, proposed in 2025, and records status \texttt{open}.
\item Supplied source attribution: Elayavalli, \emph{50 Open Problems: Ultraproduct \(\mathrm{II}_1\) Factors} (2025), Problem 17, \url{https://arxiv.org/abs/2511.20377}. The background records access on 2026-10-16 and extraction method \texttt{arxiv\_latex}; these are supplied provenance details, not an independent source visit or verification.
\item The statement contains the annotation ``Solved by J. Peterson'' and the reference key \texttt{peterson2026ultraproductapproximationspropertyt}. No expanded bibliographic citation, solution statement, or link for that reference is supplied. The attribution and solved label are reproduced as source annotations and have not been independently verified.
\end{itemize}
\end{document}
